%% Roteiro falado da apresentação: um bloco por slide.
%% Compilar: pdflatex roteiro.tex  (ou: make roteiro.pdf)
%% O ensaio com voz (ensaio_voz/ensaio.py) lê este arquivo:
%% \parte{nome}{pág. inicial}{pág. final}{tempo} e \slide{página}{título}.
%% Fala simples: como foi modelado ou o que o resultado mostra, poucos números.
\documentclass[12pt,a4paper]{article}
\usepackage[utf8]{inputenc}
\usepackage[T1]{fontenc}
\usepackage{lmodern}
\usepackage[brazil]{babel}
\usepackage[margin=2.5cm]{geometry}
\usepackage{setspace}
\usepackage{amsmath,amssymb}
\usepackage{booktabs}
\usepackage{xcolor}
\onehalfspacing
\setlength{\parindent}{0pt}
\setlength{\parskip}{0.6em}

\newcommand{\parte}[4]{\section*{#1 {\normalfont\small (páginas #2 a #3, cerca de #4)}}}
\newcommand{\slide}[2]{\par\medskip\noindent{\large\textbf{[#1] #2}}\par\nopagebreak}
\newcommand{\avanca}{\par{\hfill\small\color{gray}$\blacktriangleright$ avançar}\par}

\begin{document}

{\Large\textbf{Roteiro da apresentação}}

Um bloco por slide; o número entre colchetes é a página do PDF da
apresentação. Tempos medidos com voz sintética (\texttt{ensaio\_voz}): cerca de 29,5
minutos no total. Tempo e prioridade de cada slide em
\texttt{tempos\_prioridade.pdf}.


\parte{Abertura e Introdução}{1}{6}{2,7 min}

\slide{1}{Capa}

Bom dia a todos. Meu nome é Vitor Dall'Olivo Gonçalves e apresento o meu
Trabalho de Conclusão de Curso, ``Desenvolvimento de Sistema de Aquisição com
Estimação por Modo Deslizante para Bancada de Propulsão com Motor BLDC'',
orientado pela professora Ana Teruko Yokomizo Watanabe e coorientado por Aureo
Guilherme Dobrikopf.

\avanca

\slide{2}{Sumário}

A apresentação segue a ordem do trabalho escrito, começando pela introdução,
passando pela fundamentação teórica e pela metodologia, depois os resultados
com a discussão e, por fim, a conclusão.

\avanca

\slide{3}{Introdução}

\avanca

\slide{4}{Albatroz Aerodesign}

A motivação deste trabalho vem da equipe Albatroz Aerodesign, um projeto de
extensão da UDESC aqui de Joinville. A equipe projeta e constrói aeronaves
radiocontroladas para a competição SAE Brasil Aerodesign. Uma escolha ruim do
conjunto motopropulsor pode fazer a aeronave carregar menos peso que o
esperado, então projetar bem esse conjunto faz diferença. Desde 2022 o
regulamento só permite motores elétricos sem escovas, os chamados motores
BLDC.

\avanca

\slide{5}{Problema: ensaios fragmentados}

Para escolher esse conjunto, a equipe faz ensaios de desempenho numa bancada
como esta. Nela cada grandeza é medida de forma separada. Para saber a
rotação do motor, por exemplo, é preciso um medidor externo, à parte da
medição de tração. E o torque não é medido de jeito nenhum, mesmo sendo uma
das variáveis que mais influenciam a escolha do motor. Hoje esse valor vem do
catálogo, que foi levantado em condições diferentes das da equipe.

\avanca

\slide{6}{Objetivos}

O objetivo geral é desenvolver um sistema de aquisição de dados para ensaios de
conjuntos motopropulsores com motores BLDC.

Para chegar nisso, o trabalho se divide em cinco objetivos específicos. Medir
a tensão e a corrente nas fases do motor. Modelar o motor e projetar um
observador que estime a rotação e o torque. Reunir medição, acionamento e
segurança numa placa de medição. Ligar essa placa ao computador pela rede. E
desenvolver o programa de supervisão.

\avanca

\parte{Fundamentação teórica}{7}{12}{3,6 min}

\slide{7}{Fundamentação teórica}

\avanca

\slide{8}{Modelo do motor BLDC}

Para estimar o que acontece no eixo sem medir o eixo, primeiro é preciso
descrever o motor por equações, ou seja, modelá-lo. O modelo é dividido em
duas partes, uma elétrica e uma mecânica.

Na parte elétrica, o motor é modelado como uma resistência, uma indutância e
uma força contraeletromotriz, a FEM, em cada fase, ligadas em estrela. A FEM
é a tensão que o próprio motor gera ao girar. Como é um motor BLDC,
considera-se que ela tem forma trapezoidal e amplitude proporcional à
rotação, então ela carrega a rotação e a posição do rotor. Na parte mecânica,
o torque elétrico vem dessa FEM junto com a corrente, e o eixo é modelado com
a inércia, o atrito e a carga.

\avanca

\slide{9}{Transformadas de Clarke e Park}

A fim de reduzir o número de equações e o acoplamento entre elas, as
grandezas das três fases são levadas para dois eixos. A transformada de
Clarke passa as três fases para dois eixos fixos, alfa e beta. A
transformada de Park gira esses eixos junto com o rotor, e nesse referencial
as grandezas ficam praticamente constantes. O observador trabalha em alfa e
beta, e o torque sai da projeção das correntes medidas no ângulo do rotor.

\avanca

\slide{10}{Conjugado da hélice}

A carga no eixo do motor é a hélice, que também precisa de um modelo. Ela é
modelada com um torque parabólico, que cresce com o quadrado da rotação e é
ponderado por um coeficiente de torque, próprio de cada hélice. Esse
coeficiente depende da razão de avanço, que é a relação entre a velocidade do
ar e a rotação vezes o diâmetro da hélice.

\avanca

\slide{11}{Ponto de operação}

O ponto de operação do motor e da hélice depende da combinação dos dois, isto
é, a rotação máxima do conjunto é definida pelo cruzamento dos envelopes de
torque do motor e da hélice. É esse cruzamento que define a rotação, a
corrente e a tração do conjunto, e o torque é justamente o que a equipe não
media. As hélices usadas no trabalho são a APC 15x6E e a 17x10E.

\avanca

\slide{12}{Observador por modos deslizantes (SMO)}

A fim de estimar a rotação e o torque a partir da tensão e da corrente,
definiu-se o observador por modos deslizantes, por ser o mais empregado
comercialmente em motores BLDC. Como todo observador, ele roda um modelo das
variáveis em conjunto com o motor real e compara os dois. No modo deslizante
existe um termo chaveado, que empurra esse erro para zero. Esse sinal de
chaveamento, depois de filtrado, é justamente a FEM, e dessa estimativa saem o
ângulo e a rotação do motor.

\avanca

\parte{Metodologia}{13}{32}{10,2 min}

\slide{13}{Metodologia}

\avanca

\slide{14}{Planejamento}

O planejamento foi quebrado em frentes paralelas, a parte mecânica, a placa
de medição e a interface, para deixar o desenvolvimento mais organizado e
confortável. Essas frentes passaram por três fases. Na primeira foram
tomadas as decisões de projeto, com as restrições e definições de cada parte.
Na segunda, cada funcionalidade foi validada individualmente, primeiro em
simulação e depois no \textit{hardware}. E na terceira, tudo foi juntado e o
projeto finalizado.

\avanca

\slide{15}{Planejamento: CAD}

Na parte mecânica, primeiro foram definidos a alimentação e os equipamentos
principais. Depois, onde ficariam as placas e o esboço da bancada. E por fim
a montagem de tudo no gabinete.

\avanca

\slide{16}{Planejamento: placa de medição}

O planejamento da placa de medição foi quebrado em três frentes,
\textit{hardware}, \textit{firmware} e observador. Assim cada parte podia
avançar e ser testada sozinha. O \textit{hardware} foi do esquemático ao
protótipo, o \textit{firmware} foi validado camada por camada, e o observador
saiu da simulação para dentro da placa. No fim, as três frentes se juntam na
placa de medição completa.

\avanca

\slide{17}{Planejamento: interface de usuário}

Na interface, primeiro vieram os arquivos de configuração e as telas. Depois,
os testes de comunicação com a placa e o retorno da equipe. E por fim a
interface funcionando.

\avanca

\slide{18}{Arquitetura do sistema de aquisição}

Com o planejamento definido, o sistema foi dividido em duas partes. A placa
de medição, que mede, roda o observador e cuida da segurança. E o computador,
com a interface, os gráficos, o registro em CSV e um terminal de comandos. As
duas conversam pela rede, com o computador mandando comandos e a placa
devolvendo os dados dez vezes por segundo.

\avanca

\slide{19}{Cadeia de aquisição}

Para ensaiar o conjunto, definiu-se que o sistema precisa ligar e desligar a
bateria com segurança, medir as fases do motor, comandar o controlador e
enviar tudo ao computador. Por isso a bateria alimenta o controlador do motor
através de relés, e não diretamente.

As responsabilidades ficam divididas entre duas placas. A placa de
condicionamento, a AQ1, fica entre o controlador e o motor e só adequa as
tensões e as correntes das fases para o conversor analógico-digital. A placa
de medição concentra o resto: lê esses sinais, roda o observador, comanda os
relés e o controlador do motor.

A ligação com o computador é por Ethernet, através de um \textit{switch}. A
equipe precisava de um meio de comunicação simples e padronizado: quase todo
computador tem porta de rede cabeada, o que dispensa conversores seriais, e o
\textit{switch} permite acrescentar outras placas no futuro.

\avanca

\slide{20}{Gabinete: potência e segurança}

A parte de potência e segurança fica no gabinete. A bateria é ligada em dois
passos: primeiro um relé de pré-carga, com um resistor em série, carrega os
capacitores do controlador e evita o pico de corrente na partida; depois o
relé de potência fecha. A parada de emergência é física. E é a própria placa
que comanda o controlador do motor.

\avanca

\slide{21}{Placa de medição: camadas do firmware}

Para executar essas tarefas, a placa usa um microcontrolador RISC-V com o
sistema operacional Zephyr. O programa foi organizado em camadas. A mais
baixa cuida do acesso ao \textit{hardware} e da comunicação, a do meio reúne
os módulos do sistema, como o observador e a segurança, e a mais alta é a
aplicação. Tensões e correntes são medidas em pares, no mesmo instante, e o
observador dá um passo a cada medida.

\avanca

\slide{22}{Escolha do sensor de corrente}

Para a escolha do tipo de sensor de corrente, foram definidos os seguintes
critérios.
\begin{itemize}
  \item Isolação galvânica.
  \item Medir corrente contínua.
  \item Baixo custo, já que são três fases.
\end{itemize}
Com esses critérios, foram comparadas as seguintes opções.
\begin{itemize}
  \item \textit{Shunt} com amplificador isolado, que atende, mas é caro.
  \item Transformador de corrente e bobina de Rogowski, que não medem
        corrente contínua.
  \item Sensor Hall com núcleo toroidal, que atende aos três critérios e foi
        o escolhido.
\end{itemize}

\avanca

\slide{23}{Sensor Hall com núcleo toroidal}

No sensor escolhido, o fio da fase passa por dentro de um núcleo toroidal,
que concentra o campo magnético numa pequena fenda. O sensor Hall fica nessa
fenda e mede um campo proporcional à corrente.

\avanca

\slide{24}{Ensaio de caracterização do conjunto núcleo–sensor}

Para caracterizar o sensor, uma fonte passa uma corrente conhecida pelo
núcleo, e um multímetro lê a tensão de saída do sensor. O ensaio foi repetido
com uma, quatro e oito voltas do fio, para varrer uma faixa maior de campo.

\avanca

\slide{25}{Placa de condicionamento das fases (AQ1)}

Os sensores e as tensões das fases chegam à placa de medição através da
placa de condicionamento, a AQ1. Ela recebe os sensores de corrente e as
tensões das fases e entrega esses sinais, já filtrados, à placa de medição.

\avanca

\slide{26}{Observador embarcado}

Para levar o observador para dentro da placa, o modo deslizante descrito na
fundamentação foi modificado em dois pontos. O modo deslizante continua
estimando a FEM. Mas o ângulo do rotor não é necessariamente zero no início,
então foi adicionado um PLL, que acompanha a FEM e tira dela o ângulo e a
rotação a partir de qualquer posição. Além disso, a FEM gerada pelo modo
deslizante tem muita trepidação, uma oscilação rápida causada pelo próprio
chaveamento. Por isso a FEM é reconstruída de forma analítica a partir do
ângulo estimado, o que remove essa trepidação. Para o torque, usa-se a
corrente medida junto com o ângulo estimado.

\avanca

\slide{27}{Aplicação de supervisão}

A operação do sistema é feita pela aplicação de supervisão, no computador.
Pela tela principal se acompanha e comanda o ensaio, se ajusta o observador e
se grava tudo em CSV.

\avanca

\slide{28}{Aplicação de supervisão: calibração dos sensores}

A calibração dos sensores é feita nesta tela. Para cada canal se define o
ganho, o deslocamento de zero e os limites válidos. O valor editado vai pela
rede até a placa, que aplica o ajuste e o grava em memória não volátil. A
confirmação vem pelo próximo pacote de dados, quando a leitura do canal passa
a refletir o novo valor.

\avanca

\slide{29}{Protocolo em arquivo único}

Para que a placa e o computador se entendam, os dois precisam concordar no
formato das mensagens. Esse formato é descrito num único arquivo. Dele saem
tanto o código da placa quanto o da aplicação, então os dois lados sempre
falam a mesma língua.

\avanca

\slide{30}{Identificação dos parâmetros do motor}

O modelo do motor da fundamentação depende de parâmetros que o fabricante
não fornece por completo, então eles foram medidos. A resistência foi medida
pelo método de Kelvin, a quatro terminais. A indutância, por dois métodos
diferentes: com o motor parado, pelo método da ressonância LC, que liga
capacitores conhecidos ao enrolamento e mede a frequência de ressonância; e
com o motor girando, pelo ajuste harmônico, que ajusta o modelo aos registros
de tensão e corrente. E as perdas
mecânicas, girando o motor sem carga em várias rotações, modeladas como
atrito viscoso mais uma parcela constante.

\avanca

\slide{31}{Ajuste automático dos ganhos do observador}

Com o modelo do motor definido, falta escolher os ganhos do observador, que
foram ajustados de forma automática, em simulação. O motor passa por vários
patamares de rotação, e cada conjunto de ganhos recebe uma nota pelo erro de
torque e de rotação, com o torque pesando mais. A evolução diferencial, um
método de otimização por população de tentativas, procura a menor nota.

\avanca

\slide{32}{Ensaio do coeficiente de torque da hélice}

Por fim, o ensaio que é a finalidade da bancada: levantar o coeficiente de
torque da hélice. Ele é feito por patamares de acelerador. Em cada um,
espera-se a rotação estabilizar e registram-se rotação e torque. O torque da
hélice é o torque estimado do motor menos as perdas sem carga, e com ele se
calcula o coeficiente de torque, que é comparado com a curva do fabricante.

\avanca

\parte{Resultados e discussão}{33}{62}{10,9 min}

\slide{33}{Resultados e discussão}

\avanca

\slide{34}{Parâmetros elétricos do motor}

Os primeiros resultados são os parâmetros do motor, levantados pelos métodos
da metodologia. A resistência saiu bem consistente. Na indutância os dois
métodos não concordam: a ressonância LC, com o motor parado, dá cerca de 7
microhenries, e o ajuste harmônico, com ele girando, cerca de 11. No
observador foi usado o valor do ajuste harmônico, e essa diferença volta na discussão.

\avanca

\slide{35}{Sensor de corrente: montagem do ensaio}

Em seguida, o sensor de corrente. O ensaio de caracterização foi montado na
bancada do laboratório, como mostra a foto.

\avanca

\slide{36}{Sensor de corrente: tensão de saída $\times$ corrente}

Variando a corrente na fonte e lendo a tensão do sensor, a tensão de saída
cresce em linha reta com a corrente nas três configurações, e quanto mais
voltas de fio, mais inclinada fica a reta.

\avanca

\slide{37}{Sensor de corrente: curva $\Delta$V $\times$ B}

Como o catálogo do sensor informa a sensibilidade em campo magnético, a
corrente foi convertida em campo. Assim, as três configurações se juntam numa
mesma reta, colada na sensibilidade de catálogo e dentro da faixa do
fabricante. O sensor se mantém linear bem além da faixa em que é usado.

\avanca

\slide{38}{Simulação: desempenho do PLL}

Antes do motor real, o observador foi testado em simulação, começando pelo
PLL, que extrai o ângulo do rotor a partir da FEM estimada. O ângulo estimado
acompanha o ângulo real, e o erro entre os dois fica em poucos graus. Esse
erro é maior em rotação baixa, onde a FEM é pequena, um efeito próprio do
modo deslizante, que perde o sincronismo em valores baixos de rotação.

\avanca

\slide{39}{Ajuste automático: convergência da otimização}

Os ganhos usados nessa simulação vêm da busca automática descrita na
metodologia. Cada barra clara é uma tentativa, e a linha escura é a melhor
encontrada até aquele ponto. A convergência é relativamente rápida: em poucas
iterações a busca já converge.

\avanca

\slide{40}{Ajuste automático: ganhos obtidos por hélice}

Foi rodado o ajuste para ambas as hélices. Os ganhos do modo deslizante quase
não mudam de uma hélice para a outra, e quem muda é o ganho ligado à FEM, que
acompanha a carga. Nas duas, o erro final ficou cerca de um quinto menor que
o do ponto de partida, sem nenhum ajuste manual.

\avanca

\slide{41}{Ajuste automático: sensibilidade da função-custo a cada ganho}

Para cada um dos ganhos, rodou-se a sensibilidade em torno do ponto ótimo,
variando um ganho por vez com os demais fixos. O primeiro está no fundo do
vale, como esperado. O segundo está numa região plana, antes de uma subida
forte em que o observador perde o sincronismo. E o terceiro quase não muda o
resultado.

\avanca

\slide{42}{Simulação: rotação e torque, hélice APC 15x6E}

Com esses ganhos, simulou-se o conjunto com a hélice 15x6E, em patamares de
rotação. A rotação estimada segue a real em todos os patamares. O torque
oscila por causa do acionamento do motor, mas a média estimada fica
praticamente sobre a real.

\avanca

\slide{43}{Simulação: rotação e torque, hélice APC 17x10E}

Na 17x10E a rotação também é acompanhada. O torque médio estimado fica um
pouco acima do real, e essa diferença diminui conforme a rotação sobe.

\avanca

\slide{44}{Simulação: correntes e FEM nos eixos $\alpha\beta$}

Dentro do observador, as correntes e a FEM nos eixos alfa e beta mostram que
a estimativa segue o valor real, com a oscilação típica do modo deslizante em
volta.

\avanca

\slide{45}{Simulação: identificação do coeficiente $k_Q$ por patamar}

Por fim, o ensaio de hélice da metodologia foi refeito na simulação, onde o
coeficiente verdadeiro é conhecido. Em todos os patamares o coeficiente
estimado fica praticamente igual ao verdadeiro, com diferença abaixo de meio
por cento. Com isso, o observador se mostra teoricamente válido para estimar
os coeficientes da hélice.

\avanca

\slide{46}{Ensaio em laboratório com osciloscópio}

Validado em simulação, o observador precisava ser testado com dados reais. O
primeiro teste com o motor girando foi no laboratório, com um osciloscópio.
Foram registradas as tensões e as correntes das fases, estas já pelos
sensores da placa, além do comando do acelerador e de um tacômetro óptico,
que serve de referência de rotação.

\avanca

\slide{47}{Laboratório: perdas a vazio}

Nessas condições, com o motor girando sem hélice, foram levantadas as perdas
mecânicas em várias rotações, com o torque calculado pelo balanço de
potência. A reta se ajusta bem aos pontos, com R² de 0,93. E o tacômetro
concordou com a rotação calculada pelas tensões, com desvio abaixo de
0,1 \%.

\avanca

\slide{48}{Laboratório: observador nos registros do osciloscópio}

Aplicou-se então o observador sobre esses registros, a fim de comparar o
torque que vem do balanço de potência com o do observador. A rotação estimada
segue a do tacômetro com erro abaixo de 0,1 \%, e o torque ficou alguns por
cento abaixo da referência.

\avanca

\slide{49}{Sistema de aquisição montado}

Com o sensor e o observador validados, o sistema foi montado. A parte de
potência, o gabinete pronto e o motor com a hélice na mesa. Daqui para frente
o torque vem do observador rodando dentro da placa.

\avanca

\slide{50}{Ensaio na interface: hélice APC 15x6E}

Com o sistema montado, as duas hélices foram ensaiadas pela interface. Na
15x6E, a cada degrau de acelerador, a rotação e o torque estimados sobem
junto.

\avanca

\slide{51}{Ensaio na interface: hélice APC 17x10E}

E o mesmo para a 17x10E.

\avanca

\slide{52}{1º ensaio: coeficiente de torque das hélices}

Desses ensaios sai o coeficiente de torque, calculado como na metodologia. No
primeiro ensaio, a linha é a curva do fabricante e os pontos são as medidas.
Nas duas hélices os pontos ficaram bem acima da curva, e mais ainda em
rotação baixa.

\avanca

\slide{53}{1º ensaio: torque das hélices}

A tabela mostra o tamanho disso, de 20 a quase 50 \% acima do catálogo, um
valor acima do esperado. Além disso, ao se ajustar uma curva aos pontos, o
torque seguia a rotação elevada a um coeficiente menor do que dois, e não ao
quadrado, como se espera de uma hélice. Isso levou a considerar que houvesse
um erro nas perdas descontadas do torque.

\avanca

\slide{54}{Perdas a vazio: 1º e 2º ensaio}

Por isso as perdas foram medidas de novo, agora com o próprio sistema, na
mesma montagem do ensaio com hélice. A linha azul é a nova medida, e a
tracejada vermelha é a que tinha sido usada. As novas perdas são cerca do
dobro, e com elas as duas hélices foram ensaiadas de novo.

\avanca

\slide{55}{2º ensaio: coeficiente de torque das hélices}

No segundo ensaio, nos pontos cheios, as duas hélices passam a acompanhar a
curva do fabricante.

\avanca

\slide{56}{2º ensaio: hélice APC 15x6E}

Na 15x6E a diferença para o fabricante ficou abaixo de 3 \%.

\avanca

\slide{57}{2º ensaio: hélice APC 17x10E}

Na 17x10E ficou entre 5 e 9 \% acima. Os pontos de rotação mais baixa
dependem de uma extrapolação das perdas e merecem mais cuidado.

\avanca

\slide{58}{Discussão: o que mudou entre os ensaios}

Entre os dois ensaios, a rotação e o torque medido se repetiram, e o que
mudou foram as perdas descontadas. Afastar a mesa de objetos próximos não
mudou nada, então o vento em volta da hélice não era a causa.

A diferença nas perdas tem três causas possíveis. O método, já que no
laboratório o torque veio do balanço de potência e no segundo ensaio veio do
observador. A montagem, com outro controlador e outros cabos, cuja queda de
tensão entra na conta como se fosse do motor. E o aquecimento do motor, que
muda a resistência e o ímã, enquanto os parâmetros foram medidos a frio.

\avanca

\slide{59}{Discussão: interpretação do 2º ensaio}

Quando as perdas e o torque com hélice vêm do mesmo observador, na mesma
montagem, um eventual erro do observador aparece nos dois e se cancela em boa
parte. Então o resultado mostra que as perdas devem ser medidas nas mesmas
condições do ensaio.

\avanca

\slide{60}{Discussão: indutância, corrente e montagem}

Sobre a indutância, o método com o motor parado trabalha numa frequência bem
mais alta e com corrente muito menor que na operação, e o método com o motor
girando inclui os cabos. Um ensaio com corrente de operação resolveria isso.
A corrente fica limitada pelo sensor a cerca de 60 amperes, o que bastou para
estas hélices.

\avanca

\slide{61}{Discussão: memória do microcontrolador}

Este gráfico mostra como a memória do microcontrolador ficou ocupada. A maior
parte vai para a rede e para o sistema operacional, e sobrou pouco espaço
livre. Por isso a placa não distribui endereço de rede sozinha, e o
computador precisa ser configurado à mão.

\avanca

\slide{62}{Limitações}

Resumindo as limitações: não há uma medida de torque independente no eixo, a
indutância ficou incerta, a corrente tem um teto, a memória impede a conexão
automática, e a tração e o vento não foram medidos.

\avanca

\parte{Conclusão}{63}{68}{2,0 min}

\slide{63}{Conclusão}

\avanca

\slide{64}{Objetivos atingidos}

Retomando os objetivos, o sistema mede as tensões e correntes do motor, o
motor foi caracterizado e o observador estima a rotação com erro abaixo de
0,1 \%. A placa reúne medição, observador, acionamento e segurança, e com ela
as hélices ficaram próximas da curva do fabricante. A placa conversa com o
computador pela rede, e a aplicação comanda e registra os ensaios.

\avanca

\slide{65}{Contribuições e trabalhos futuros}

Como contribuições ficam uma placa de baixo custo que estima rotação e torque
sem sensor no eixo, os ensaios definidos no computador sem precisar regravar
a placa, e o ajuste automático do observador, que pode ser repetido para
outros motores.

Como trabalhos futuros, ficam os seguintes pontos.
\begin{itemize}
  \item Um dinamômetro, para ter uma medida de torque de referência direto no
        eixo e comparar com o observador.
  \item Células de carga para medir a tração, um tubo de Pitot para a
        velocidade do ar e ensaios em túnel de vento, saindo da condição
        estática.
  \item Levar o observador para a região de enfraquecimento de campo, onde o
        motor passa da rotação nominal.
  \item Estudar a bateria, levantando a relação entre o estado de carga e a
        tensão em aberto, e o efeito da temperatura.
  \item Uma bancada física com seis graus de liberdade, para ensaiar o
        conjunto em movimento.
\end{itemize}

\avanca

\slide{66 e 67}{Referências}

\textit{(Não ler. Passar direto.)}

\avanca\avanca

\slide{68}{Obrigado}

Agradeço a atenção e fico à disposição para as perguntas da banca.

\end{document}
